% =============================================================================
% 06_DISCUSSION — Interpretation, Implications, Limitations, Agenda
% =============================================================================
% Page target: ~2 pages
% Structure: Summary → Implications → Limitations → Research Agenda
% =============================================================================

\section{Discussion}
\label{sec:discussion}
\addcontentsline{toc}{section}{Discussion}

This discussion interprets the synthesized findings presented in
Section~\ref{sec:results} through an integrated analytical lens, examines
their implications for research, policy, and clinical practice, acknowledges
the limitations of this review, and concludes with a forward-looking
research agenda identifying the most pressing unanswered questions at the
intersection of health AI evaluation and governance.

\subsection{Summary of Key Findings}
\label{sec:discussion_summary}

The three thematic findings presented in the results are analytically
interdependent rather than independent. The \textbf{ecological validity
gap} in evaluation methodology does not merely reflect a technical
measurement problem — it has direct consequences for governance. When AI
systems are validated under conditions that do not reflect real-world
clinical complexity, governance frameworks based on those validations
will be structurally misaligned with the phenomena they seek to govern.
Similarly, the \textbf{governance capacity deficit} documented in health
systems is not simply a matter of underinvestment — it is partly a
consequence of evaluation frameworks that generate overconfident
validation evidence (suggesting AI systems are more ready for deployment
than they are), which reduces the perceived urgency of building
pre-deployment governance infrastructure. Finally, the \textbf{pilot
trap} can be understood as the downstream consequence of these two
upstream failures: systems validated on inadequate evidence, governed
by frameworks that assume absent capacity, predictably fail to transition
to routine care.

This interdependency implies that improving health AI outcomes requires
simultaneous action across all three dimensions — better evaluation
methodology, governance capacity investment, and implementation support
— rather than sequential approaches that treat evaluation, governance,
and implementation as separate phases.

\subsection{Implications}
\label{sec:discussion_implications}

\subsubsection{Implications for AI Evaluation Research}

The reviewed literature makes clear that the field needs a paradigm shift
from single-site, retrospective evaluation to prospective, multicenter,
ecologically grounded validation. This requires: (i) embedding AI evaluation
within real clinical workflows rather than controlled experimental setups;
(ii) adopting longitudinal evaluation designs that capture algorithmic drift
and performance evolution over time; (iii) developing clinical harm-asymmetric
performance metrics that reflect the differential consequences of false-positive
and false-negative errors in specific clinical contexts; and (iv) mandating
subgroup performance reporting to detect and quantify algorithmic bias across
population strata. The TRAC-Pain protocol \citep{bailes2026trac} offers a
methodological template that merits wider adoption and adaptation.

\subsubsection{Implications for Regulatory and Policy Practice}

The regulatory implications cut in two directions. First, current regulatory
frameworks must evolve from one-time pre-market authorization toward
\textbf{lifecycle-spanning governance} that mandates continuous post-market
performance monitoring, requires disclosure of training data sources, and
establishes enforceable equity and bias mitigation standards
\citep{abulibdeh2025illusion, warraich2024fda}. Second, regulators must
recognize that governance capacity in health systems is not pre-given and
cannot be assumed — it must be actively built through implementation
support, not just prescribed through standards \citep{tian2026pilottrap}.
This suggests that regulatory frameworks should incorporate implementation
support mechanisms, including regulatory sandboxes, technical assistance
programs, and mandatory pre-deployment organizational readiness assessments.

\subsubsection{Implications for Health Systems and Clinical Practice}

For health systems, the reviewed evidence points to the need to invest
in organizational governance infrastructure as a precondition for AI
deployment, not as an afterthought. This includes: establishing clear
accountability structures for AI-assisted clinical decisions; building
cross-functional governance committees that bring together clinical,
technical, legal, and patient representatives; developing operational
protocols for model updates, performance monitoring, and incident
response; and investing in workforce development to build algorithmic
literacy and clinical-AI workflow integration skills at scale.

\subsection{Limitations}
\label{sec:discussion_limitations}
